% Generado por generar-historial.sh el 2026-09-24. No editar a mano.
\begin{table}[H]
\caption{Commits del Trabajo Parcial por integrante y entregable}\label{tab:commits-autor}
\footnotesize
\begin{tabular}{@{}lrrr@{}}
\toprule
\textbf{Integrante} & \textbf{PC1 (hasta \texttt{pc1})} & \textbf{PC2 (nuevos)} & \textbf{Merges de PR (PC2)} \\
\midrule
Dayana Kety Gómez Rodriguez & 5 & 0 & 0 \\
Rody Sebastian Vilchez Marin & 21 & 22 & 15 \\
\bottomrule
\end{tabular}

{\footnotesize\textit{Nota.} Commits que tocan \texttt{tp/} o \texttt{docs/} en todas las ramas publicadas en \texttt{origin}. PC1 cuenta hasta el tag \texttt{pc1} (versión entregada el 13 de septiembre); PC2, lo que no está en ese tag. Los merges de PR se cuentan aparte porque no tocan archivos por sí mismos.\par}
\end{table}

{\small
\begin{xltabular}{\textwidth}{@{}L{1.9cm}L{2.7cm}L{1.4cm}Yc@{}}
\caption{Commits de la PC2, del más reciente al más antiguo}\label{tab:commits}\\
\toprule
\textbf{Fecha} & \textbf{Autor} & \textbf{Commit} & \textbf{Mensaje} & \textbf{En \texttt{main}} \\
\midrule
\endfirsthead
\toprule
\textbf{Fecha} & \textbf{Autor} & \textbf{Commit} & \textbf{Mensaje} & \textbf{En \texttt{main}} \\
\midrule
\endhead
\bottomrule
\endfoot
2026-09-24 & Rody Sebastian Vilchez Marin & \texttt{4199e93} & Medir la laptop con el protocolo completo y sumarla a la comparación & no \\
2026-09-24 & Rody Sebastian Vilchez Marin & \texttt{3e717da} & Nombrar los entornos por lo que son en la tabla de resultados & no \\
2026-09-23 & Rody Sebastian Vilchez Marin & \texttt{9755b44} & Correr el K-means en un Pixel 9a y medirlo con el protocolo de la PC2 & no \\
2026-09-23 & Rody Sebastian Vilchez Marin & \texttt{0f0d7a3} & Contrastar el K-means en GPU contra la CPU en la misma máquina & no \\
2026-09-23 & Rody Sebastian Vilchez Marin & \texttt{42b4a66} & Mantenimiento: notebook sin salidas, ignorar artefactos de Spin por semana y checkout v5 & no \\
2026-09-23 & Rody Sebastian Vilchez Marin & \texttt{b4d62b0} & Marcar como fusionadas las PRs 9, 21 y 22 en la tabla del informe & no \\
2026-09-23 & Rody Sebastian Vilchez Marin & \texttt{9f256a0} & Atender la revisión de Sourcery: tamaños imposibles, k del archivo y el invariante de Datos & no \\
2026-09-23 & Rody Sebastian Vilchez Marin & \texttt{00734d0} & Regenerar el historial de commits con la rama del informe publicada & no \\
2026-09-23 & Rody Sebastian Vilchez Marin & \texttt{53b2bf6} & Registrar la PR del informe en la tabla de pull requests & no \\
2026-09-23 & Rody Sebastian Vilchez Marin & \texttt{0167830} & Redactar el informe de la PC2 en LaTeX con tablas generadas desde los reportes & no \\
2026-09-17 & Rody Sebastian Vilchez Marin & \texttt{c92ed2f} & Cerrar el análisis con el barrido de chunk y la reproducibilidad exacta & no \\
2026-09-17 & Rody Sebastian Vilchez Marin & \texttt{0d62fa9} & Analizar speedup, escalabilidad y recursos con los datos medidos & no \\
2026-09-17 & Rody Sebastian Vilchez Marin & \texttt{bb48eb4} & Recolectar memoria y CPU por configuración en un solo comando & no \\
2026-09-17 & Rody Sebastian Vilchez Marin & \texttt{a2969f1} & Medir memoria y CPU por fase, y generar las tablas desde el JSON & no \\
2026-09-17 & Rody Sebastian Vilchez Marin & \texttt{4d34600} & Verificar la sincronización del K-means con Spin & no \\
2026-09-17 & Rody Sebastian Vilchez Marin & \texttt{20a1707} & Medir el speedup con un protocolo fijado, no improvisado & no \\
2026-09-17 & Rody Sebastian Vilchez Marin & \texttt{82ca48b} & Agregar la CLI del K-means con metadatos de reproducibilidad & no \\
2026-09-15 & Rody Sebastian Vilchez Marin & \texttt{b1bd4bc} & Implementar el K-means concurrente con worker pool persistente & no \\
2026-09-15 & Rody Sebastian Vilchez Marin & \texttt{e027bea} & Implementar el K-means secuencial y la inicialización k-means++ & no \\
2026-09-15 & Rody Sebastian Vilchez Marin & \texttt{62d73fe} & Leer el gold en un arreglo plano de N*6 float64 & no \\
2026-09-15 & Rody Sebastian Vilchez Marin & \texttt{393f3fd} & Fijar el diseño del K-means concurrente antes de implementarlo & no \\
2026-09-12 & Rody Sebastian Vilchez Marin & \texttt{37e66a6} & Verificar en CI que los notebooks no lleguen con salidas & no \\
\end{xltabular}
}

{\footnotesize\textit{Nota.} «En \texttt{main}» indica si el commit ya forma parte de \texttt{origin/main} al generar esta tabla; los que dicen «no» están en una rama publicada pendiente de fusión.\par}
